% Part: conditional-logics
% Chapter: minimal-change-semantics
% Section: true-false

\documentclass[../../../include/open-logic-section]{subfiles}

\begin{document}

\olfileid[tr]{con}{min}{tf}

\olsection{Karşıolgusalların Doğruluğu ve Yanlışlığı}

Bir karşıolgusal $!A\cif !B$, \olref{fig:true}'da gösterildiği gibi en
yakın $!A$-dünyalarının tümü $!B$-dünyalarıysa (boşuna olmayan biçimde)
doğrudur.
\begin{figure}
\begin{center}
\begin{tikzpicture}[scale=.7]\small
  \spheresystem{5}
  \draw (0,0) node {$w$};
  \propositionintersect{0}{4}{40}{3.5}
  {\tikzset{proposition/.append={smooth,tension=1.4}}
  \proposition[shift={(1.6,-1)}]{90}{1}{30}{4}}
  \path (1.5,3) node[below] {$\formula{B}$};
  \path (3,0) node {$\formula{A}$};
\end{tikzpicture}
\caption{Boşuna olmayan biçimde doğru karşıolgusal}
\ollabel{fig:true}
\end{center}
\end{figure}
$w$ çevresindeki küreler sisteminde $!A$'yı kabul eden hiçbir küre yoksa
karşıolgusal $w$'de yine doğrudur. Bu durumda \emph{boşuna} doğrudur
(bkz. \olref{fig:vacuous}).
\begin{figure}
\begin{center}
\begin{tikzpicture}[scale=.7]\small
  \spheresystem{5}
  \draw (0,0) node {$w$};
  \proposition{0}{6.5}{40}{3.5}
  {\tikzset{proposition/.append={smooth,tension=1.4}}
  \proposition[shift={(1.2,-1)}]{90}{1}{30}{4}}
  \path (1.5,3) node[below] {$\formula{B}$};
  \spherepos{0}{7}{node {$\formula{A}$}}
\end{tikzpicture}
\caption{Boşuna doğru karşıolgusal}
\ollabel{fig:vacuous}
\end{center}
\end{figure}

İki biçimde yanlış olabilir. Birincisi, en yakın $!A$-dünyalarının tümü
$!B$-dünyaları değilken bazılarının $!B$-dünyası olmasıdır. Bu durumda
$!A\cif\lnot !B$ de yanlıştır (bkz. \olref{fig:false}).
\begin{figure}
\begin{center}
\begin{tikzpicture}[scale=.7]\small
  \spheresystem{5}
  \draw (0,0) node {$w$};
  \propositionintersect{0}{4}{40}{3.5}
  {\tikzset{proposition/.append={smooth,tension=1.4}}
  \proposition[shift={(1.6,0)}]{90}{1}{30}{3}}
  \path (1.5,3) node[below] {$\formula{B}$};
  \path (3,0) node {$\formula{A}$};
\end{tikzpicture}
\caption{Yanlış karşıolgusal, yanlış karşıtı}
\ollabel{fig:false}
\end{center}
\end{figure}
En yakın $!A$-dünyaları $!B$-dünyalarıyla hiç örtüşmüyorsa
$!A\cif !B$ yanlıştır. Fakat bu durumda en yakın bütün $!A$-dünyaları
$\lnot !B$-dünyalarıdır; dolayısıyla $!A\cif\lnot !B$ doğrudur
(bkz. \olref{fig:false-opposite}).
\begin{figure}
\begin{center}
\begin{tikzpicture}[scale=.7]\small
  \spheresystem{5}
  \draw (0,0) node {$w$};
  \propositionintersect{0}{4}{40}{3.5}
  {\tikzset{proposition/.append={smooth,tension=1.4}}
  \proposition[shift={(-1,0)}]{90}{1}{30}{3}}
  \path (-1,3) node[below] {$\formula{B}$};
  \path (3,0) node {$\formula{A}$};
\end{tikzpicture}
\caption{Yanlış karşıolgusal, doğru karşıtı}
\ollabel{fig:false-opposite}
\end{center}
\end{figure}

Sıkı koşuldan farklı olarak karşıolgusallar olumsal olabilir.
\olref{fig:contingent}'deki küre modelini düşünün. $u$'ya en yakın
$!A$-dünyalarının tümü $!B$-dünyaları olduğundan
$\mSat{M}{!A\cif !B}[u]$. Fakat $v$'ye en yakın $!A$-dünyalarının bazıları
$!B$-dünyası değildir; dolayısıyla $\mSat/{M}{!A\cif !B}[v]$.
\begin{figure}
\begin{center}
\begin{tikzpicture}[scale=.6]\small
  \spheresystem{7}
  \spheresystem[shift={(0:3.1)},dashed]{4}
  \propositionintersect{45}{4}{40}{4.5}
  \begin{scope}
    \clip \propositionplot{45}{4}{40}{4.5} ;
    \spherelayer[shift={(0:3.1)}]{3}
  \end{scope}
  \proposition[shift={(1.4,-.5)}]{90}{1}{35}{5}
  \path (0,0) node {$u$};
  \path (0:3.1) node {$v$};
  \spherepos{35}{9}{node {$\formula{A}$}}
  \spherepos{80}{9}{node {$\formula{B}$}}
\end{tikzpicture}
\end{center}
\caption{Olumsal karşıolgusal}
\ollabel{fig:contingent}
\end{figure}
\end{document}
